% Shared preamble of the HSEMHS mode documents (docs/physics/*.tex, docs/technical/*.tex).
\documentclass[10pt]{extarticle}
\usepackage[a4paper,margin=1.7cm]{geometry}
\usepackage[T1]{fontenc}\usepackage{lmodern}\usepackage{tabularx}
\usepackage{amsmath,amssymb,bm,enumitem,titlesec,xcolor,booktabs,array,tikz,graphicx,longtable}
\usepackage[hidelinks]{hyperref}
\titleformat{\section}{\large\bfseries\color{blue!50!black}}{\thesection.}{0.5em}{}
\titlespacing*{\section}{0pt}{9pt}{3pt}
\titleformat{\subsection}{\normalsize\bfseries}{\thesubsection}{0.5em}{}
\titlespacing*{\subsection}{0pt}{6pt}{2pt}
\setlist{nosep,leftmargin=1.3em}
\setlength{\parindent}{0pt}\setlength{\parskip}{4pt}
\renewcommand{\arraystretch}{1.15}
\newcommand{\B}{\mathbf B}\newcommand{\J}{\mathbf J}\newcommand{\F}{\mathbf F}\newcommand{\R}{\mathbf R}
\newcommand{\zh}{\hat{\mathbf z}}\newcommand{\Bo}{\mathbf B_{\rm obs}}\newcommand{\Lv}{\mathbf L}
\newcommand{\code}[1]{\texttt{#1}}
\newcolumntype{L}[1]{>{\raggedright\arraybackslash}p{#1}}
\newcolumntype{Y}{>{\raggedright\arraybackslash}X}
\newcommand{\doctitle}[2]{{\Large\bfseries #1}\\[2pt]{\small #2}\par\vspace{2pt}}
\pagestyle{plain}
\AtBeginDocument{\sloppy}

\begin{document}
\doctitle{mhs-poisson-z: pressure-only magnetohydrostatic correction (height grid)}{Physics of \code{run.mode = "mhs-poisson-z"} (HSEMHS). 24 Oct 2026.}

\section{What the mode does and why}
The full magnetohydrostatic solve (\code{full-mhs-z}) changes both the pressure and the magnetic field and needs a day
on four nodes for the Bifrost cube. Inside a Stokes inversion, where the atmosphere changes at every iteration, that is
far too slow. The FIRTEZ-dz inversion code (Borrero et al.) instead keeps $T$ and $\B$ as given and corrects
\textbf{only the pressure}, by solving a Poisson-type equation. This mode does the same on our grid with our equation of
state: it starts from \code{column-hse-z}, keeps $T$ and $\B$ fixed, and finds the pressure that balances gravity and
the given Lorentz force as well as possible. It takes minutes.

\section{Quantities}
\begin{tabular}{>{$}l<{$} p{11.0cm} l}
\toprule
\text{Symbol} & Meaning & Unit \\ \midrule
T,\ \B & temperature and magnetic field from the input, \textbf{both fixed} & K, T \\
\J,\ \Lv & current density $\nabla\times\B/\mu_0$; Lorentz force $\Lv=\J\times\B$ (fixed) & A\,m$^{-2}$, N\,m$^{-3}$ \\
P,\ u & gas pressure, the only unknown; solved as $u=\ln P$ & Pa, -- \\
\rho & $\rho(T,P)$ from the EOS table & kg\,m$^{-3}$ \\
H,\ \bar H & scale height $P/(\rho g)$; its layer median in the starting state & m \\
\beta & plasma beta $2\mu_0P/B^2$ & -- \\
z_{\rm cut} & top of the solved region (production 1\,Mm) & m \\
\bottomrule
\end{tabular}

\section{The equation}
A static atmosphere has zero net force per unit volume,
\[
\F\equiv-\nabla P+\Lv-\rho g\,\zh=0 .
\]
With $\B$ fixed, $\Lv$ is a given force field, and $\rho$ depends on $P$ only through the EOS. Taking the divergence,
\[
\nabla^2P=\nabla\cdot\Lv-g\,\partial_z\rho ,
\]
a Poisson equation for $P$; hence the name. In a horizontally periodic box it separates into Fourier modes, each a
one-dimensional problem in $z$, which is what makes it fast.

\section{What a pressure correction can and cannot do}
Any force field splits uniquely into a \textbf{gradient part} $-\nabla\phi$ (curl-free) and a \textbf{non-gradient part}
$\mathbf S$ (divergence-free, orthogonal to every gradient), with $\int|\F|^2=\int|\nabla\phi|^2+\int|\mathbf S|^2$.
A pressure change adds a pure gradient, $-\nabla\delta P$. If $\rho$ were fixed, it could remove the gradient part
exactly and never touch $\mathbf S$.

Which forces are non-gradient? The curl of the force balance is
\[
\nabla\times\F=\nabla\times\Lv-g\,(\partial_y\rho,\,-\partial_x\rho,\,0),
\]
so static balance needs the \textbf{curl condition} $\nabla\times(\J\times\B)=g(\partial_y\rho,-\partial_x\rho,0)$: the
magnetic tension must hold up every horizontal density contrast. Gravity is a gradient only if $\rho$ is horizontally
uniform; buoyancy of hot granules next to cool lanes is non-gradient. Without a field the only static atmosphere is
horizontally uniform. With Bifrost's $T$ and fixed $\B$ the curl condition does not hold, so \emph{no pressure can make
the atmosphere exactly static}.

\textbf{The density loophole.} $\rho$ is not fixed: with $T$ fixed, $\delta\rho=\rho\,(\partial\ln\rho/\partial\ln P)_T\,\delta P/P$,
and $\delta(\nabla\times\F)=-g\nabla\times(\delta\rho\,\zh)$. By changing the pressure, and hence the density, the
correction can shrink horizontal density contrasts and part of $\mathbf S$, but not the part due to the temperature
contrast at fixed pressure or to the tension of $\B$. Measured on the production cube (original script, 23 Oct): column
HSE is 80--90\,\% gradient error in the photosphere; the Poisson correction cut the gradient part 3.5--5 times and,
through the density, the non-gradient part by 1.4--3.2 between 0 and 0.75\,Mm.

\section{The quantity minimised}
The force is divided by $P$ and the unknown is $u=\ln P$:
\[
\R\equiv\frac{\F}{P}=-\nabla u-\frac{\zh}{H(T,P)}+\frac{\Lv}{P},\qquad
\Phi(u)=\sum_{\text{centres}}w\,\bar H^2\big(R_x^2+R_y^2\big)+\sum_{z\text{-faces}}h\,\bar H^2R_z^2 .
\]
\begin{itemize}
\item $P=e^u$ stays positive across the many decades in the box.
\item $\R$ is a \emph{relative} imbalance: minimising $|\F|^2$ itself rewards lowering $P$ (the first version drove $P$
toward zero in 637\,000 cells).
\item $\R\bar H\approx\F/(\rho g)$: the imbalance is measured in units of the weight of the gas, layer by layer, so the
dense deep layers do not swamp the upper ones. $w,h$ are the cell heights.
\item Without $\Lv$, $R_z=0$ is $du/dz=-1/H$, i.e.\ column HSE.
\end{itemize}
The mean of $\ln P$ on the bottom plane is held at its column-HSE value: the overall pressure level is set by the input,
not by the correction.

\section{The low-$\beta$ problem and $z_{\rm cut}$}
Where $\beta\ll1$ the term $\Lv/P$ dominates. Static balance there requires a force-free field, which a snapshot of a
moving simulation is not; with $\B$ fixed, the only way to balance a large Lorentz force is a large pressure gradient.
Because $\partial P/\partial z$ couples the layers, this inflates the pressure in the photosphere too: solved up to
3\,Mm the median $\ln P$ rose by 3.2 (a factor 25); with the Lorentz force switched off nothing happened. The force
balance is therefore solved only below $z_{\rm cut}=1$\,Mm (median $\beta\approx8$); above it each column is integrated
hydrostatically from the solved pressure at $z_{\rm cut}$. FIRTEZ makes the same choice in a softer way
(\code{LIMIT\_MAGTERM}). Above $z_{\rm cut}$ the atmosphere is hydrostatic, \textbf{not} magnetohydrostatic.

\section{Optical depth}
As in \code{column-hse-z}: $\tau_{500}$ is integrated down from the top of the box with the \code{exp\_tail} closure,
and $z(\tau_{500}=1)$ is reported per column.

\section{Results on the production cube}
Mode run of 24 Oct (8 ranks on 4 nodes, 4.7\,min, column HSE included): $\ln(P/P_{\rm colHSE})$ below $z_{\rm cut}$ has
percentiles 1/50/99 $=-0.54/+0.06/+0.99$; $z(\tau_{500}=1)$ percentiles 1/50/99 $=-5/69/125$\,km against $-29/64/127$\,km
in column HSE. The objective fell to 8\,\% of its starting value; the run used the then step limit (60) while still
improving by $2\times10^{-4}$ per step. Force imbalance per layer from the original script (23 Oct, same physics), RMS
of $|\F|$ over the layer median of $\rho g$:
\begin{center}
\begin{tabular}{rcccc}
\toprule
$z$ [Mm] & column HSE & Poisson correction & Bifrost snapshot & \code{full-mhs-z} (M5h, iteration 132) \\ \midrule
$-0.5$ & 0.11 & \textbf{0.06} & 0.11 & 0.09 \\
$0.0$  & 0.33 & \textbf{0.13} & 0.29 & 0.27 \\
$0.25$ & 0.23 & \textbf{0.07} & 0.12 & 0.19 \\
$0.5$  & 0.22 & \textbf{0.06} & 0.12 & 0.18 \\
$0.75$ & 0.28 & \textbf{0.09} & 0.21 & 0.22 \\
$1.0$  & 0.66 & \textbf{0.34} & 0.49 & 0.51 \\
\bottomrule
\end{tabular}
\end{center}
This diagnostic uses the same discretisation the Poisson step minimises, which favours it; \code{full-mhs-z} balances
forces on its own staggered grid.

\section{What it means}
\begin{itemize}
\item Where photospheric lines form ($z\lesssim0.5$\,Mm) the correction removes most of the force imbalance of column
HSE, below that of the Bifrost snapshot itself, and barely moves the $\tau$ scale.
\item The remaining imbalance is the buoyancy of the fixed temperature structure that the fixed field cannot hold up.
Removing it needs $\B$ (or $T$) to change: \code{full-mhs-z}.
\item $z_{\rm cut}$ is a physical choice: where the fixed field stops being trusted to be in equilibrium.
\end{itemize}
\end{document}
